\chapter{Animations: spinners library}
\label{ch:spinners}

\begin{aside}
Ten spinners. The same loop. The choice is which one fits
your aesthetic.
\end{aside}

\section{Catalogue}

ASCII (4 frames): \code{| / - \textbackslash}

Dots (3 frames): \code{. .. ...}

Arrow (4 frames): \code{\^{} > v <}

Earth (3 frames): three globe glyphs.

Moon (8 frames): phases from new to full.

Braille (10 frames): rotating dot pattern using Unicode
braille.

Box (4 frames): rotating outline.

Pong (8 frames): a ball bouncing along a track.

Spiral (4 frames): nested arcs.

Triangle (3 frames): rotating triangle.

\section{Choosing}

Match the spinner to the app's tone:

\begin{itemize}[leftmargin=*]
  \item Serious tools: ASCII or dots.
  \item Playful tools: pong, earth, moon.
  \item Modern: braille (smooth).
\end{itemize}

\section{Tick rate}

50--150 ms per frame. Faster feels frantic; slower feels
sluggish.

\section{Multiple spinners}

If many spinners are visible, they share a tick. All move
together; the wall clock drives the frame index.

\section{Stop conditions}

A spinner stops when its driving condition (loading
signal) clears. The frame index resets so the next time
it appears, it starts at frame 0.

\section{Recap}

\begin{recap}
\begin{itemize}[leftmargin=*]
  \item Many spinners, one loop.
  \item Match tone to app.
  \item 50--150 ms per frame.
  \item Shared tick across spinners.
\end{itemize}
\end{recap}

\section*{Workshop}

\begin{enumerate}[leftmargin=*]
  \item Try each spinner; pick the one you like.
  \item Build a custom 6-frame spinner.
  \item Share a tick across three concurrent spinners.
\end{enumerate}
